% rmq.tex -- private working draft of the RMQ manuscript substrate.
% Branch codex/eg-cp-paper-evidence-r1, base commit
% 745a3c5b1fe99fe1072fbcf03f18f5dc9c23cb82.
% This draft asserts only theorem surfaces present at that base commit.
% The packed all-size architecture result is provisional and appears only at
% the single marked insertion point in Section 9.
\documentclass[11pt]{article}

\usepackage[T1]{fontenc}
\usepackage[utf8]{inputenc}
\usepackage{amsmath}
\usepackage{amssymb}
\usepackage{amsthm}
\usepackage[margin=1.1in]{geometry}
\usepackage{booktabs}
\usepackage{hyperref}

% Ledger anchors: \ledger{<id>} ties the preceding mathematical claim to
% its row in paper/THEOREM_LEDGER.md. The macro renders nothing; the checker
% paper/check_paper.ps1 enforces bidirectional coverage.
\newcommand{\ledger}[1]{}

% Lean identifiers.
\newcommand{\leanid}[1]{\texttt{#1}}

\newtheorem{theorem}{Theorem}[section]
\newtheorem{lemma}[theorem]{Lemma}
\theoremstyle{definition}
\newtheorem{definition}[theorem]{Definition}
\newtheorem{targetstmt}[theorem]{Target Statement}
\theoremstyle{remark}
\newtheorem{remark}[theorem]{Remark}

\newcommand{\flog}[1]{\lfloor\log_2 #1\rfloor}
\newcommand{\len}[1]{\lvert #1\rvert}

\title{Payload-Accounted Succinct Range Minimum Queries,\\
Machine-Checked in Lean~4:\\
Reference Semantics, a Charged-Trace Cost Model,\\
and an Explicit Trust Boundary\thanks{Private working draft. This document
is an internal manuscript substrate pinned to repository commit
\texttt{745a3c5b1fe99fe1072fbcf03f18f5dc9c23cb82}. It is not a submission,
not a public claim surface, and not an architecture-acceptance record.}}
\author{RMQ repository contributors}
\date{Draft of 2026-08-05}

\begin{document}
\maketitle

\begin{abstract}
We describe a Mathlib-free Lean~4 development that mechanizes a succinct
range-minimum-query (RMQ) family over ordinary integer lists together with
the accounting needed to state its space and query cost without ambiguity.
The development proves: (i) a payload bound
$\len{\mathrm{buildPayload}(xs)} \le 2n + \mathrm{overhead}(n)$ with a
checked little-$o$-linear overhead, where $n$ is the input length;
(ii) exact leftmost half-open RMQ answers for every valid query and uniform
rejection of invalid queries; (iii) a uniform modeled query-cost bound: the
canonical execution emits an explicit word-read trace whose length is at most
the checked constant $210$, derived from a named component algebra rather
than asserted; and (iv) an information-theoretic lower bound for exact RMQ
encodings, stated in an integer-arithmetic ``doubled Catalan slack'' form
equivalent to $2n - 1.5\log_2 n - O(1)$ bits. A model-adequacy layer ties
the cost claim to an explicit trace and store: every charged event is an
attempted word read, no synthetic cost-only markers occur, successful reads
are backed by counted payload words, the physical word list erases exactly
to the public payload, and store-parametricity holds under agreement on the
execution's ordered read footprint, with checked mutation witnesses showing
that the store is genuinely observed. We state precisely what is and is not
charged, keep payload bits, allocated bits, charged probes, model ticks, and
Lean wall-clock time distinct, and document the trust base (Lean kernel
checking under a pinned toolchain, standard axioms only). A packed
cell-probe strengthening, bounding the allocated capacity of one packed
object and the probes into it, has been accepted in the repository but is
not yet absorbed into this draft; we reserve exactly one marked insertion
point for it and assert it nowhere else. Novelty statements are deliberately conditional pending a
completed search log.
\end{abstract}

\section{Introduction}\label{sec:intro}

Range minimum query is a canonical static data-structure problem: given an
array, preprocess it so that the position of the minimum in a queried
subrange can be returned quickly. The classical succinct answer stores
$2n + o(n)$ bits and answers queries in constant time in a word-RAM or
cell-probe model~\cite{FischerHeun11,Fischer10}, and the space is essentially
optimal because the answer function determines the Cartesian tree shape of
the input, of which there are Catalan-many~\cite{Vuillemin80,GBT84}.

Mechanizing this story in a proof assistant raises a difficulty that is
easy to underestimate: the informal literature freely mixes several
quantities that a formal artifact must keep apart --- the bits of the
advertised payload, the bits actually allocated by a concrete memory object,
the number of memory probes charged to a query, the ticks of whatever
abstract cost model the development uses, and the wall-clock behavior of the
proof assistant's executable fragment. A formalization that silently
conflates any two of these can prove a theorem that reads like a succinct
data-structure result while actually asserting something weaker. The
repository this draft describes was built around that failure mode: its
theorem surfaces name their model explicitly, its cost claims are traces
rather than aggregate constants, and its documentation separates non-claims
from claims.

\subsection{Contributions of the mechanization}\label{sec:contributions}

All results below are kernel-checked Lean~4 theorems at repository commit
\texttt{745a3c5b1fe99fe1072fbcf03f18f5dc9c23cb82}; the ledger
(\texttt{paper/THEOREM\_LEDGER.md}) maps every claim in this draft to an
exact declaration and source file.

\begin{enumerate}
\item \emph{Reference semantics.} A value-level half-open, leftmost-tie RMQ
  contract over \leanid{List Int}, with uniqueness of the specified answer
  (Section~\ref{sec:semantics}).
\item \emph{Payload-accounted succinct upper bound.} A construction whose
  advertised payload is at most $2n + \mathrm{overhead}(n)$ bits with
  $\mathrm{overhead}$ checked little-$o$-linear, exact answers on valid
  queries, uniform rejection of invalid queries, and a uniform modeled cost
  bound with the checked value $210$
  (Section~\ref{sec:upper}).\ledger{L-UB-01}
\item \emph{An explicit charged-trace execution model.} The costed query is
  the projection of an explicit event trace; every emitted event is an
  attempted word read; the trace is backed by one concrete physical word
  list that erases to the public payload; store-parametricity and
  answer dependency on the store are checked, including a decisive
  single-word corruption witness (Section~\ref{sec:model-adequacy}).
\item \emph{An information-theoretic lower bound.} Every exact fixed-length
  RMQ state encoding needs $2n - 1.5\log_2 n - O(1)$ payload bits, proved in
  an integer-arithmetic doubled form, with a concrete $2n$-bit upper witness
  in the same encoding class (Section~\ref{sec:lower}).\ledger{L-LB-01}
\item \emph{Reusable spokes.} Standalone rank/select families (a plain
  Jacobson/Clark-style $n + o(n)$ family and a fixed-weight compressed/FID
  family) and the balanced-parentheses close-navigation layer consumed by
  the RMQ capstone (Section~\ref{sec:upper}).
\end{enumerate}

\subsection{What this draft does not claim}\label{sec:nonclaims-intro}

The development makes no claim about compiled Lean execution speed, compiler
correctness, CPU or memory-hierarchy semantics, or production serialization.
The modeled cost bound charges attempted payload-word reads only; controller
dispatch, decoding, arithmetic, comparisons, and branching are uncharged in
the current model, so the constant $210$ is a charged-trace bound in an
explicitly named model, not a running-time bound in the conventional
word-RAM model (Section~\ref{sec:quantities}). The space theorem of
Section~\ref{sec:upper} counts advertised payload bits, not the allocated
capacity of a packed memory object (Section~\ref{sec:quantities}); the
allocated-capacity strengthening is a separate accepted result over a
different object, which this draft states in one place and nowhere else
(Section~\ref{sec:pending}). Priority claims are deliberately conditional:
closely related mechanizations exist, most prominently the
Coq/SSReflect succinct-data-structure line of Tanaka, Affeldt, Garrigue,
and Qi~\cite{TAG16,AGQT19}, and the search log
(\texttt{paper/RELATED\_WORK\_LEDGER.md}) is not yet complete
(Section~\ref{sec:related}).

\subsection{Draft status}\label{sec:status}

This is a private working draft assembled from the repository's accepted
claim maps at one exact commit. One result family --- a packed
counted-header cell-probe refinement --- is under a separate feasibility
gate and is neither accepted nor rejected at the base commit. This draft
reserves a single insertion point for it (Section~\ref{sec:pending}) and
otherwise excludes it from every claim.

\section{Reference semantics}\label{sec:semantics}

The specification layer is deliberately value-level and list-based: the
reference objects are ordinary Lean lists of unbounded integers, and the
reference answer is defined by arithmetic on list lookups, not by any data
structure. All structure-level theorems refine this one contract.

\begin{definition}[Valid query]\label{def:validrange}
For $xs : \leanid{List Int}$ and naturals $\mathit{left}, \mathit{right}$,
the query $[\mathit{left}, \mathit{right})$ is \emph{valid}, written
$\leanid{ValidRange}\ xs\ \mathit{left}\ \mathit{right}$, when
$\mathit{left} < \mathit{right} \wedge \mathit{right} \le \len{xs}$.
A valid query is a nonempty half-open window inside the list.
\end{definition}

\begin{definition}[Leftmost argmin]\label{def:leftmostargmin}
$\leanid{LeftmostArgMin}\ xs\ \mathit{left}\ \mathit{right}\ \mathit{idx}$
holds when the query is valid, $\mathit{left} \le \mathit{idx} <
\mathit{right}$, and there is a value $v$ with $xs[\mathit{idx}] = v$ such
that $v \le w$ for every value $w$ at a position in
$[\mathit{left}, \mathit{right})$, and $v < w$ strictly for every value $w$
at a position in $[\mathit{left}, \mathit{idx})$. Ties therefore resolve to
the leftmost minimizing index.
\end{definition}

\begin{lemma}[Uniqueness]\label{lem:unique}
For fixed $xs$, $\mathit{left}$, $\mathit{right}$, at most one index
satisfies \leanid{LeftmostArgMin} \textup{(}Lean:
\leanid{leftmostArgMin\_unique}\textup{)}.\ledger{L-REF-01}
\end{lemma}

The executable reference implementation is a linear window scan
(\leanid{scanWindow}), proved to return the \leanid{LeftmostArgMin} witness
on every valid window; several intermediate backends (sparse table, hybrid
block decompositions, Fischer--Heun-style structures) are each proved
extensionally equal to this reference on their domains. The succinct
construction of Section~\ref{sec:upper} is compared against the same
contract, so ``correct'' means one fixed thing across the repository.

A second load-bearing specification fact is \emph{shape sufficiency}: the
answer to every RMQ query is determined by the Cartesian tree shape of the
input~\cite{Vuillemin80,GBT84}, so a representation may discard values and
retain shape. The repository formalizes Cartesian shapes, their
balanced-parentheses (BP) encodings, and the RMQ/LCA reductions over
Euler-tour depth sequences with the $\pm 1$ depth
invariant~\cite{HarelTarjan84,BFC00}, and proves that the succinct query
over the shape representative agrees with the value-level reference
answer.\ledger{L-REF-02}

\section{Quantities, models, and the accounting discipline}
\label{sec:quantities}

Five quantities are kept distinct throughout the development and throughout
this draft. Conflating any two of them changes the meaning of a succinct
data-structure claim, so we fix vocabulary once, here.

\begin{description}
\item[Payload bits.] The length of the advertised bit-level payload
  $\leanid{buildPayload}(xs) : \leanid{List Bool}$. This is what the
  $2n + o(n)$ upper bound of Section~\ref{sec:upper} counts. Proof-only
  fields, certificates, and Lean-structure overhead are excluded by
  construction, and the repository's anti-vacuity lints reject theorem
  shapes that would smuggle reference computation into an uncounted field.
\item[Allocated bits.] The full capacity of a concrete cell array,
  including any header, final-cell padding, sentinel or dead cells. This is
  strictly stronger than payload bits, and neither theorem of
  Section~\ref{sec:upper} establishes it: the $2n + o(n)$ bound there counts
  the advertised payload, and the physical-erasure theorem of
  Section~\ref{sec:model-adequacy} equates payload \emph{bit contents} with
  the flattened word list. Allocated capacity is bounded separately, and by
  a different object --- the packed cell-probe representation of
  Section~\ref{sec:pending}, whose complete capacity, counting the header
  cell, every payload cell, and final padding at full cell width, is at most
  $2n + \rho(n)$ for a checked little-$o$-linear
  $\rho$.\ledger{L-PACK-01}
\item[Charged probes (model reads).] Events of the explicit trace layer.
  On the accepted route every charged event is one attempted read of one
  word of the modeled store, including failed reads. The constant $210$
  bounds the number of such events per query.
\item[Model ticks.] The \leanid{Costed} cost of a query. On the canonical
  route this is proved \emph{equal} to the emitted trace length, so ticks
  and charged probes coincide there by theorem, not by definition.
\item[Lean wall-clock time.] The observable behavior of Lean's evaluator on
  the executable fragment. No theorem in the repository mentions it; the
  validation harness measures it separately and it plays no role in any
  claim of this draft.
\end{description}

\paragraph{The charge policy.}
The declared policy of the current model is: \emph{charged} --- attempted
payload-word reads (\leanid{readWord}), one model tick each, including
failed reads; \emph{uncharged} --- instruction dispatch, register moves,
fixed-width decoding of a fetched word, bounded arithmetic and comparisons
on register values, option tests, branching, candidate merging, trace
assembly, and the validity guard. After the charged-fringe and same-block
recharges, every uncharged step on the accepted route is a
bounded-per-step register computation between two charged reads; no loop
remains on the accepted route whose trip count grows with the input and
that is silent in the trace. The chunk caps making this checkable are
described in Section~\ref{sec:upper}.\ledger{L-UB-19}

\paragraph{What the model is not.}
Because controller work is uncharged, the bound $210$ is not a conventional
word-RAM running-time result, and this draft never presents it as one. It
is closer in spirit to a cell-probe accounting --- memory touches are the
unit --- but the current theorem does not package a cell-probe object
either: the store is word-addressed by logical segment, not a single packed
bit array with a frozen width function. That packaging is precisely what the
accepted packed cell-probe result of Section~\ref{sec:pending} supplies, over
a different object than the theorem discussed here. Charging an
instruction-level controller is a separate planned strengthening
(Section~\ref{sec:limitations}).\ledger{L-OPEN-06}

\section{The information-theoretic lower bound}\label{sec:lower}

The number of Cartesian shapes of size $n$ is the Catalan number $C_n$, and
distinct shapes induce distinct RMQ answer functions, so any exact encoding
must distinguish $C_n$ objects. The repository states the resulting bound
for a concrete encoding class: an \leanid{ExactRMQStateEncoding}~$n$~$b$ is
a $b$-bit-budget state encoding equipped with a decoder that answers every
valid query exactly on all representative arrays of size $n$, together with
a payload view that counts the bits its states charge.

To avoid rational arithmetic in a Mathlib-free development, the public
statement is doubled. Let
$\mathrm{dSlack}(n) = 4n - (3\flog{2n+1} + 3)$
(Lean: \leanid{doubledLogSlackLower}, with $\flog{\cdot}$ the natural
floor logarithm \leanid{Nat.log2}).

\begin{theorem}[Doubled Catalan-slack payload bound]\label{thm:lower}
For every $n$:
\begin{enumerate}
\item for every $b$ and every exact encoding
  $E : \leanid{ExactRMQStateEncoding}\ n\ b$,
  \(\mathrm{dSlack}(n) \le 2b\);
\item the same doubled lower bound applies to any uniform charged payload
  budget: if every on-domain built state of $E$ charges at most
  $\mathit{budget}$ payload bits, then
  \(\mathrm{dSlack}(n) \le 2\,\mathit{budget}\);
\item there exists an exact encoding with budget $2n$ whose built state
  charges exactly $2n$ payload bits on every representative shape of
  size~$n$.
\end{enumerate}
\textup{(}Lean:
\leanid{exactRMQ\_tight\_fixed\_length\_payload\_space\_bound\_doubled\_catalan\_slack},
public alias
\leanid{RMQ.Headlines.exactRMQLowerBoundDoubledCatalanSlack}.\textup{)}
\ledger{L-LB-01}
\end{theorem}

Halving the doubled inequality recovers the familiar reading: every exact
encoding needs at least $2n - 1.5\flog{2n+1} - 1.5$ bits, i.e.\
$2n - 1.5\log_2 n - O(1)$, and the third clause shows the $2n$-bit side is
attained by a concrete decoder rather than a bare counting existential. An
undoubled variant with the weaker slack
$2n - (2\flog{2n+1} + 2)$ is also checked
(Lean: \leanid{exactRMQ\_tight\_fixed\_length\_payload\_space\_bound}).
\ledger{L-LB-02}

\paragraph{Scope.}
This is a payload-capacity counting bound for exact encodings. It is
\emph{not} the cell-probe lower bound of Liu and Yu~\cite{LiuYu20}, later
sharpened by Liu~\cite{Liu21}, which bounds the redundancy of \emph{any}
structure achieving query time $t$ in the cell-probe model. Those results
are related work only; nothing of them is mechanized in this repository,
and this draft cites them purely for context.\ledger{L-OPEN-05}

\section{The succinct upper bound}\label{sec:upper}

\subsection{Construction architecture}\label{sec:architecture}

The construction follows the classical shape-based recipe: encode the
Cartesian shape of the input as a balanced-parentheses string of exactly
$2n$ bits, and attach $o(n)$ bits of auxiliary directories so that queries
can navigate it with a constant number of charged reads. Concretely, the
advertised payload $\leanid{buildPayload}(xs)$ splits into:

\begin{itemize}
\item the BP shape code ($2n$ bits);
\item rank/access directories in the Jacobson/Clark
  style~\cite{Jacobson89,Clark96}: sampled superblock and block counters
  over the BP code;
\item a sparse-exception select directory for close-parenthesis selection;
\item four-Russians-style chunk tables for fringe windows: precomputed
  per-chunk excess/minimum summaries consumed by the endpoint-fringe and
  same-block legs, in the spirit of the microtable decompositions of the
  classical RMQ literature~\cite{FischerHeun11};
\item a compact close/LCA directory over a relative range-min-max
  macrostructure, playing the role that range-min-max trees play in
  succinct tree navigation~\cite{NavarroSadakane14,MunroRaman01};
\item an interior directory with sparse sampled levels for cross-block
  interior minima.
\end{itemize}

The query pipeline maps a valid window $[\mathit{left},\mathit{right})$ to
close positions via select, runs the close/LCA leg (endpoint fringes by
chunked reads, interior by the sparse directory), and maps the winning
close position back to an index via rank.

\subsection{Space}\label{sec:space}

\begin{theorem}[Payload bound]\label{thm:payload}
For every $xs : \leanid{List Int}$ with $n = \len{xs}$,
\[
  \len{\leanid{buildPayload}(xs)} \;\le\; 2n + \mathrm{overhead}(n),
\]
where $\mathrm{overhead}$ is a single closed function of $n$.
\textup{(}Lean: \leanid{buildPayload\_length}.\textup{)}\ledger{L-UB-01}
\end{theorem}

\begin{theorem}[Little-$o$ overhead]\label{thm:littleo}
$\mathrm{overhead}$ is little-$o$-linear in the repository's
integer-arithmetic sense: for every positive scale $s$ there is a threshold
$N$ such that $s \cdot \mathrm{overhead}(n) \le n$ for all $n \ge N$.
\textup{(}Lean: \leanid{overhead\_littleO}, predicate
\leanid{SuccinctSpace.LittleOLinear}.\textup{)}\ledger{L-UB-02}
\end{theorem}

No padding is inserted to force an exact size equality; the bound is an
honest inequality about the constructed payload. We emphasize again that
this counts payload bits, not allocated capacity
(Section~\ref{sec:quantities}), and that the overhead envelope is proved
little-$o$ but not proved tight: no theorem exhibits a reachable input
family attaining the envelope (Section~\ref{sec:limitations}).
\ledger{L-OPEN-04}

\subsection{Correctness}\label{sec:correctness}

\begin{theorem}[Exact answers on valid queries]\label{thm:exact}
For every $xs$ and every valid query $[\mathit{left},\mathit{right})$, the
succinct query returns \leanid{some}~$i$ where $i$ is the unique
\leanid{LeftmostArgMin} index of Definition~\ref{def:leftmostargmin};
moreover the same exactness transfers to any caller-supplied store that
agrees with the canonical store on the checked read footprint.
\textup{(}Lean:
\leanid{listIntFinalFullModelSoundnessExactOfFootprintGlobal}; packaged in
the paper-facing
\leanid{RMQ.Headlines.listIntSuccinctRMQPaperMainTheorem}.\textup{)}
\ledger{L-UB-03}
\end{theorem}

\begin{theorem}[Invalid-query boundary]\label{thm:invalid}
Empty, reversed, and out-of-bounds ranges return \leanid{none}, uniformly
across the plain, costed, trace, and supplied-store query surfaces.
\textup{(}Lean: \leanid{queryCosted\_invalid},
\leanid{queryCosted\_empty\_range}, \leanid{queryCosted\_reversed\_range},
\leanid{queryCosted\_out\_of\_bounds}.\textup{)}\ledger{L-UB-04}
\end{theorem}

\subsection{Modeled query cost}\label{sec:cost}

\begin{theorem}[Uniform charged-trace cost bound]\label{thm:cost}
Every query has modeled cost at most \leanid{queryCost}, and
$\leanid{queryCost} = 210$ by a checked equality derived from the component
algebra, uniformly for all input sizes.
\textup{(}Lean: \leanid{queryCost\_eq};
\leanid{concreteBPNativeSuccinctRMQPrincipledAllSizeChargedTraceCost\_eq}.\textup{)}
\ledger{L-UB-05}
\end{theorem}

\begin{theorem}[Component algebra]\label{thm:algebra}
The constant decomposes as
\[
  2\cdot\underbrace{35}_{\textup{select}}
  \;+\; \bigl(2\cdot\underbrace{11}_{\textup{rank}}
  + 2\cdot\underbrace{37}_{\textup{fringe}}
  + \underbrace{33}_{\textup{interior}}\bigr)
  \;+\; \underbrace{11}_{\textup{rank}} \;=\; 210,
\]
i.e.\ two select legs, a close/LCA leg (two rank seeds, two endpoint-fringe
windows, one interior directory pass), and one final rank. The equality is
derived from named per-component caps, not asserted as a literal.
\textup{(}Lean:
\leanid{chargedTraceCostAlgebra}.\textup{)}\ledger{L-UB-06}
\end{theorem}

Two structural theorems make the constant meaningful rather than
decorative:

\begin{theorem}[Read-only event vocabulary]\label{thm:readword}
Every event emitted by the canonical whole-query trace is a
\leanid{readWord} event --- an attempted read of one word of the modeled
store. \textup{(}Lean:
\leanid{concreteBPNativeSuccinctRMQWholeQueryGlobalWordTraceResult\_readWord\_only}.\textup{)}
\ledger{L-UB-07}
\end{theorem}

\begin{theorem}[No synthetic cost-only events]\label{thm:nosynthetic}
The dedicated synthetic cost-only trace constructor, used historically to
migrate aggregate costs, does not occur in the canonical all-size trace.
\textup{(}Lean: the no-synthetic execution story
\leanid{succinctRMQGlobalPayloadStoreNoSyntheticExecutionStory} and its
flat-payload strengthening.\textup{)}\ledger{L-UB-08}
\end{theorem}

\begin{theorem}[Certificate weight equalities]\label{thm:weight}
On the canonical trace, the sum of per-event non-synthetic weights equals
the emitted trace length and equals the \leanid{Costed} cost of the same
execution, and is at most $210$; a counterfactual theorem shows that
inserting a synthetic event anywhere breaks the weight/length equality.
\textup{(}Lean:
\leanid{concreteBPNativeSuccinctRMQWholeQueryGlobalWordTraceResult\_nonSyntheticWeight\_sum\_le\_210}
and companions.\textup{)}\ledger{L-UB-09}
\end{theorem}

\paragraph{Bounded uncharged remainder.}
The charge policy of Section~\ref{sec:quantities} is checkable because the
formerly event-silent scans were recharged: fringe windows read at most $33$
chunk-table entries plus $4$ window words (cap $37$); at most $8$ chunks fit
one machine word; the same-block close leg is charged by the same chunk fold
with the same cap $37$, with no size hypothesis, so the caps hold at
$n \in \{0,1,2\}$ too. \textup{(}Lean:
\leanid{bpChunkedSameBlockCloseSeededCosted\_cost\_le} and the cap
identities in \leanid{ChargedFringeChunks}, \leanid{ChargedWordChunks},
\leanid{ChargedTableRegime}.\textup{)}\ledger{L-UB-19}

\begin{remark}[Why the constant moved, and why that is the policy working]
\label{rem:constant-moved}
The current literal $210$ supersedes a retired value $207$: recharging the
interior directory's sparse-level reads made three genuine reads visible in
the trace, and the literal absorbed them. The distinction is between
algorithmic work (comparisons, merges --- unchanged) and representation
artifacts (memory touches needed to extract an operand from the chosen
encoding). A memory-touch cost model must charge both; leaving a
representation artifact uncharged is not a cheaper algorithm but an
unmodelled one. The retired literal remains frozen as a checked historical
identity rather than deleted, so no later recharge can silently rewrite the
record. No claim of minimality attaches to $210$
(Section~\ref{sec:limitations}).\ledger{L-UB-06}
\end{remark}

\subsection{Reusable spokes}\label{sec:spokes}

Two standalone theorem families support the capstone and are exposed
independently. First, a plain-bitvector rank/select family in the
Jacobson/Clark style: stored-bit access, exact rank and select, counted
payload $n + \mathrm{overhead}(n)$ with little-$o$-linear overhead, and a
constant modeled query bound \textup{(}Lean:
\leanid{jacobsonClarkNPlusOConstantQuery}\textup{)}.\ledger{L-RS-01}
Second, a fixed-weight compressed/FID family in the
Raman--Raman--Rao line~\cite{RRR02}: an enumerative fixed-weight primary
payload plus $o(n)$ auxiliary payload, exact access/rank/select, and one
uniform constant modeled query bound, together with an interpreted replay
routing the charged reads through the first-order word-RAM bridge layer
\textup{(}Lean: \leanid{compressedFIDFixedWeightFamilyProfile},
\leanid{compressedFIDFixedWeightInterpretedFamilyProfile}\textup{)}.
\ledger{L-RS-02}
Third, the BP close-navigation layer consumed by the RMQ path: $o(n)$
auxiliary close-navigation payload, constant modeled cost, exact
answer-close semantics for Cartesian-shape queries with exact endpoint
close positions, and machine-word-bounded payload reads \textup{(}Lean:
\leanid{concreteBPCloseNavigationFamily\_profile}\textup{)}.\ledger{L-BP-01}

These spokes are formalized to the depth the capstone needs; the repository
does not claim a complete succinct tree-navigation library in the sense of
Navarro--Sadakane~\cite{NavarroSadakane14}.

\section{Execution-model adequacy}\label{sec:model-adequacy}

A constant like $210$ is only as meaningful as the link between the cost
model and an execution that a skeptic can inspect. The adequacy layer makes
that link explicit, and every clause below is a checked theorem at the base
commit.

\subsection{One physical word list, erasing to the public payload}

The reviewer route fixes one pre-execution physical word list, assembled
from an exhaustive typed universe of $22$ physical sources covering $23$
logical segments $0..22$ (the two BP roles share one physical source), with
exclusive regions and complete segment coverage. Its erasure --- the
flattening of its bit contents --- is exactly the public
$\leanid{buildPayload}$. \textup{(}Lean:
\leanid{concreteBPNativeSuccinctRMQReviewerPhysicalWords\_erases}.\textup{)}
\ledger{L-UB-10}
This is an equality of payload bit contents; it is deliberately \emph{not}
an allocated-cell or padded-capacity statement
(Section~\ref{sec:quantities}).

\subsection{Genuine supplied-store execution}

The public physical execution is not a replay of a precomputed trace: the
supplied-store evaluator reads the caller's flat store through a checked
address-translation adapter, and canonical flat-physical execution refines
the logical execution while preserving the decoded result, modeled cost,
ordered successful and failed reads, repeated reads, and the
execution-derived footprint. \textup{(}Lean:
\leanid{succinctRMQReviewerPhysicalExecutionRefinesLogical},
\leanid{succinctRMQReviewerPhysicalStoreAdapter}.\textup{)}\ledger{L-UB-11}

\begin{theorem}[Footprint determinacy]\label{thm:footprint}
Agreement of two stores on the first execution's consumed ordered read
footprint determines the complete physical execution --- trace, failures,
repetitions, and returned value. The list-facing corollary: a supplied
store agreeing with the canonical global store on the checked footprint
yields the canonical result, the exact RMQ answer, and the transferred cost
bound. \textup{(}Lean:
\leanid{concreteBPNativeSuccinctRMQWholeQueryFlatPhysicalTraceResultWithStore\_eq\_of\_orderedFootprint};
\leanid{queryTraceResultWithStore\_eq\_of\_orderedReadFootprint};
\leanid{listIntFinalFullModelCostLeOfFootprintGlobal}.\textup{)}
\ledger{L-UB-12}\ledger{L-UB-18}
\end{theorem}

\begin{theorem}[The store is observed]\label{thm:corruption}
Differing translated supplied-store evaluator values yield differing
physical answers, and a checked decisive single-word corruption --- changing
one consumed table word (logical segment $21$, index $3$) from the five-bit
encoding of $1$ to the five-bit encoding of $4$ --- changes the whole-query
answer from \leanid{some 0} to \leanid{none}. \textup{(}Lean:
\leanid{succinctRMQReviewerPhysicalValueDependency} and the validation
guards.\textup{)}\ledger{L-UB-13}
\end{theorem}

\begin{theorem}[Successful reads are payload-backed]\label{thm:backing}
Under footprint agreement with the canonical store, every successful read
in the supplied-store replay is backed by the counted flat payload --- the
anti-oracle clause: answers cannot be fed from proof-only fields or
uncounted certificates. \textup{(}Lean:
\leanid{...WithStore\_successful\_reads\_backed\_by\_counted\_flat\_payload\_of\_footprint\_global}.\textup{)}
\ledger{L-UB-14}
\end{theorem}

\subsection{Provenance, liveness, and mutation rejection}

For the exact current query, every indexed read in the trace retains its
global position, producing program-instruction occurrence, folded prefix
state, component-local position, exact invocation parameters, source, and a
multiplicity-preserving embedding into the composed trace.
\textup{(}Lean:
\leanid{succinctRMQReviewerEveryReadOccurrenceProvenance}.\textup{)}
\ledger{L-UB-17}
Query-independently, the manifest packet proves that every counted source
and every shared-BP consumer has a successful occurrence through some
actual closed whole-query execution under a valid list query; that the
successful-occurrence predicate implies the common mutation-side predicate;
and that a counterfactual fresh segment $23$ --- given a plausible existing
label --- fails that predicate. Source regions are exclusive, segment
coverage is complete, and legacy duplicate close/interior storage is
excluded. \textup{(}Lean:
\leanid{succinctRMQReviewerManifestSemanticAdequacy}.\textup{)}
\ledger{L-UB-15}

\subsection{Bounded event data and the machine certificate}

One query-independent width,
$\leanid{reviewerWordBits}(n) = \leanid{machineWordBits}(400000\cdot(n+1))$,
bounds every stored and returned word, translated address, segment
encoding, query operand, primitive operand/result, and consumed footprint
address, with the checked all-size inequality
$\leanid{reviewerWordBits}(n) \le 20\cdot(\flog{n+2}+1)$.
\ledger{L-UB-16}
These facts are collected in a $24$-field machine certificate whose fields
are consumed, one by one, by an independent expected-type record, and
lifted to the \leanid{List Int} surface by a guarded adequacy packet; the
paper-facing main theorem consumes the certificate's cost field by literal
projection rather than restating a numeral. \textup{(}Lean:
\leanid{succinctRMQReviewerMachineWellFormed},
\leanid{succinctRMQReviewerMachineRequiredFacts},
\leanid{listIntSuccinctRMQReviewerNativeMachineAdequacy}.\textup{)}
\ledger{L-UB-20}

\paragraph{Replayable anti-bypass evidence.}
Beyond the kernel theorems, the repository maintains a committed mutation
replay of $41$ cases expecting $40$ semantic rejections and one
packet-only expected accept, exercising shape, sibling-store, failed-probe,
disconnected-trace, and consumer bypasses. We flag its evidence tier
precisely: replays are reproducible-artifact evidence, one tier below the
kernel theorems above, and nothing in this draft rests on them alone.

\section{Mechanization and the trust boundary}\label{sec:trust}

\paragraph{Toolchain and dependencies.}
The development is Lean~4~\cite{MouraUllrich21}, pinned by
\texttt{lean-toolchain} to Lean~4.22.0, and is deliberately Mathlib-free:
it uses the Lean/Std library plus the \leanid{omega} tactic shipped with
Lean. This is an engineering choice made to keep the checked import closure
small and auditable, not a judgment about Mathlib.

\paragraph{What must be trusted.}
Theorem truth rests on Lean kernel checking of the committed declarations
under the pinned toolchain. The checked surfaces use only the standard
non-computational axioms that Lean can introduce
(\leanid{propext}, \leanid{Quot.sound}, \leanid{Classical.choice});
there are no project-specific axioms, and the repository gate rejects
\leanid{sorry}/\leanid{sorryAx}, custom \leanid{axiom} declarations,
\leanid{unsafe}, \leanid{implemented\_by}, \leanid{extern},
\leanid{native\_decide}, and \leanid{ofReduceBool} in checked source. Axiom
inventories are printed by per-surface \leanid{\#print axioms} scripts.

\paragraph{What need not be trusted.}
The validation executable, cost harness, example programs, lints, and the
mutation replay are reviewer and reproducibility checks; they are not part
of the trust base, and no theorem depends on them. Likewise, the
development process (AI-assisted, audit-driven) is recorded as provenance;
process evidence explains decisions but proves no mathematics.

\paragraph{Import discipline.}
The narrow paper-facing import root is \leanid{RMQPaper}, whose checked
closure contains the succinct RMQ construction, the lower bound, and the
word-RAM/model-adequacy machinery, and excludes legacy, compatibility,
archive, and proposal barrels. Historical cost surfaces remain
kernel-checked but are quarantined behind explicitly named
compatibility aliases that \leanid{RMQPaper} does not import.

\section{Reproducibility}\label{sec:repro}

All claims of this draft refer to repository commit
\texttt{745a3c5b1fe99fe1072fbcf03f18f5dc9c23cb82}. The artifact reproduction
path is:

\begin{itemize}
\item \texttt{bash scripts/reproduce\_artifact.sh} --- build, public-root
  builds, validation executable, headline/word-RAM/full axiom checks,
  forbidden-token and reduction-shortcut scans, whitespace checks;
\item \texttt{powershell -File scripts/gate.ps1} --- the Windows gate
  equivalent;
\item \texttt{lake env lean scripts/headline\_axiom\_check.lean} --- the
  smallest command printing the axiom inventory of the public aliases cited
  here;
\item \texttt{powershell -File paper/check\_paper.ps1} --- the
  deterministic manuscript checker for this draft (citation closure,
  duplicate labels and keys, forbidden tokens, theorem-ledger coverage, and
  the single pending-architecture insertion point);
\item \texttt{latexmk -pdf rmq.tex} (in \texttt{paper/}) --- the PDF build
  used for this draft.
\end{itemize}

Continuous integration runs the repository gate on a pinned Linux image;
validation runs report payload overhead, model cost, and wall-clock time in
separate columns, maintaining the quantity separation of
Section~\ref{sec:quantities} in the measurements as well as in the
theorems. Measurements never substitute for the space or cost theorems.

\section{Toward a packed cell-probe theorem (provisional)}
\label{sec:pending}

The strongest form of the succinct story packages the entire structure as
\emph{one} read-only array of fixed-width cells and charges the query for
every aligned cell probe, computation between probes being free --- the
standard convention in which succinct RMQ upper and lower bounds are
usually read~\cite{Fischer10,LiuYu20}. The repository's endgame roadmap
freezes such a target:

\begin{targetstmt}[Accepted in the repository; not yet absorbed into this draft]
\label{tgt:packed}
For every input list of length $n$, preprocessing constructs one read-only
array of exact-width $w(n)$-bit cells, with checked all-size bounds on $w$
yielding the conventional logarithmic-width corollary. Its complete
allocated capacity --- including a constant serialized header and final-cell
padding --- is at most $2n + \rho(n)$ bits for a checked little-$o$-linear
$\rho$. For every valid half-open query, a closed uniform controller, whose
dynamic inputs are exactly $n$, the endpoints, and the words returned by
its previous probes into that same array, returns the leftmost minimum
index after at most $C$ adaptive probes, for one exact constant $C$
independent of $n$.\ledger{L-ARCH-01}
\end{targetstmt}

At the base commit this statement is no longer a target: the feasibility
gate closed and the repository recorded a positive architecture acceptance,
discharging the geometry-closure and controller-closure rows on one exact
candidate lineage together with a fresh-blind exact-commit audit and a
coordinator reconstruction. The theorem ledger reproduces the repository's
status label accordingly.

What has \emph{not} happened is editorial: this draft has not yet absorbed
the accepted statement into a theorem environment with its proof obligations
spelled out at the level the rest of the manuscript uses. Until it does, the
marked insertion point below remains the only location licensed to carry the
result, and no other section states, assumes, or paraphrases it as a theorem
of this draft. Two consequences of the acceptance are already reflected
above, because leaving them out would have made the draft assert falsehoods:
the allocated-bits entry of Section~\ref{sec:quantities} now points here
rather than denying that any theorem bounds capacity, and the corresponding
limitation has been withdrawn.

\subsection{Result insertion point}\label{sec:insertion}

Exactly one location in this manuscript may carry the packed all-size
result. The acceptance condition has been met; the remaining step is
editorial, and until it is taken the marker below is the entire content of
that location, and no other section states, assumes, or paraphrases the
result as a theorem of this draft.

\medskip
\noindent\verb|ARCHITECTURE_RESULT_PENDING|
\medskip

\noindent The marker is retained deliberately rather than deleted: it is
what the manuscript checker keys on to confirm that the result appears in
exactly one place, and it will be replaced by the theorem statement and its
proof obligations in the same edit that removes this paragraph. Every
theorem above is already checked at the base commit and none depends on the
result carried here, so that edit is additive.

\section{Related work}\label{sec:related}

\paragraph{Classical RMQ and LCA.}
Cartesian trees originate with Vuillemin~\cite{Vuillemin80}; the
RMQ-to-$\pm 1$-RMQ and RMQ/LCA reductions with linear-space constant-time
solutions descend from Gabow--Bentley--Tarjan~\cite{GBT84},
Harel--Tarjan~\cite{HarelTarjan84}, and
Bender--Farach-Colton~\cite{BFC00}. Fischer and
Heun~\cite{FischerHeun11} gave the systematic $2n + o(n)$-bit,
constant-time succinct RMQ preprocessing scheme, and
Fischer~\cite{Fischer10} the non-systematic optimal-succinctness variant
whose model vocabulary the packed cell-probe result of
Section~\ref{sec:pending} follows.

\paragraph{Succinct building blocks.}
Rank/select dictionaries begin with Jacobson~\cite{Jacobson89} and
Clark~\cite{Clark96}; compressed indexable dictionaries with
Raman--Raman--Rao~\cite{RRR02}; balanced-parentheses navigation with
Munro--Raman~\cite{MunroRaman01} and, in fully-functional range-min-max
form, Navarro--Sadakane~\cite{NavarroSadakane14}. Navarro's
text~\cite{Navarro16} is the modern reference for the surrounding
vocabulary. The repository's spokes (Section~\ref{sec:spokes}) mechanize
the fragments of this toolkit that the RMQ capstone consumes.

\paragraph{Lower bounds.}
The information-theoretic bound mechanized here
(Section~\ref{sec:lower}) is the counting bound implicit in the classical
literature. The substantially stronger cell-probe redundancy lower bounds
for succinct RMQ are due to Liu and Yu~\cite{LiuYu20} and
Liu~\cite{Liu21}; they are not mechanized here.

\paragraph{Mechanized succinct data structures.}
The closest precedent is the Coq/SSReflect line of Tanaka, Affeldt, and
Garrigue, who verified the Jacobson rank algorithm with extraction to
OCaml~\cite{TAG16}, extended with Qi to tree algorithms over LOUDS and
related succinct representations~\cite{AGQT19}. Those papers already state
priority for mechanized succinct-structure verification in their setting;
accordingly, this draft claims no such priority. The present development
differs in prover (Lean~4, Mathlib-free), in problem focus (RMQ end-to-end
from reference semantics to a payload-accounted capstone with an
information-theoretic lower bound in the same repository), and in its
charged-trace anti-oracle discipline (Section~\ref{sec:model-adequacy});
whether any component of it is novel relative to all prior mechanizations
is left conditional on the completed search log
(\texttt{paper/RELATED\_WORK\_LEDGER.md}).

\paragraph{Mechanized cost analysis.}
Verified cost reasoning in proof assistants has several traditions:
Nipkow's amortized analyses of functional structures~\cite{Nipkow15,
Nipkow16}, time credits in separation logic with the
Chargu\'eraud--Pottier union-find verification as flagship~\cite{CP19},
asymptotic-complexity frameworks over imperative programs~\cite{GCP18,
ZhanHaslbeck18}. Those efforts tie cost to program steps of an executable
or deeply embedded program. The present development occupies a different
design point: cost is defined by an explicit model-level trace over a
declared store, with adequacy theorems (Section~\ref{sec:model-adequacy})
substituting for a program-logic connection. We present this as a
trade-off, not a superiority claim: the model is simpler and the
accounting is exhaustive over reads, but no compiled-code statement
follows. The Isabelle Archive of Formal Proofs~\cite{AFP} hosts many
verified data-structure analyses; our search of it for succinct
rank/select or RMQ entries is recorded, with its limitations, in the
related-work ledger.

\section{Limitations and open problems}\label{sec:limitations}

Stated as theorems-not-yet-proved rather than hedges; each carries a ledger
row.

\begin{enumerate}
\item \emph{Controller charging.} Controller dispatch, decoding, bounded
  arithmetic, comparisons, and branching are uncharged; a richer
  instruction-level machine that charges each such step and simulates the
  same execution is future work. Consequently no claim is made about
  running time in the conventional word-RAM model.\ledger{L-OPEN-06}
\item \emph{Preprocessing.} The construction's preprocessing complexity
  (time and workspace, in any model) is unproved and unclaimed.
  \ledger{L-OPEN-01}
\item \emph{Tightness of the constant.} $210$ is an upper cap with an
  exact reachable interior-cost witness at one component, but global
  minimality across component correlations is not proved; the constant is
  an artifact of the current representation and charge policy.
  \ledger{L-OPEN-02}
\item \emph{Tightness of the overhead envelope.} $\mathrm{overhead}$ is
  little-$o$-linear but not proved tight; no reachable family is shown to
  attain the envelope, and the envelope is not claimed comparable to
  Fischer--Heun redundancy.\ledger{L-OPEN-04}
\item \emph{Cell-probe lower bounds.} The Liu--Yu/Liu redundancy bounds
  are not mechanized; the mechanized bound is the counting bound of
  Section~\ref{sec:lower}.\ledger{L-OPEN-05}
\item \emph{Internal audit status.} All theorems cited here are
  kernel-checked at the base commit. The project's own acceptance process
  additionally requires fresh-blind audits for some surfaces; at the base
  commit the audit of the release lineage of the $210$ chain is still open
  in that internal process. This is a process status, not a mathematical
  caveat, and it is recorded here for candor.
\end{enumerate}

\section{Conclusion}\label{sec:conclusion}

The development connects, in one kernel-checked repository: a value-level
RMQ reference contract; Cartesian-shape and BP reductions; a
payload-accounted $2n + o(n)$ succinct construction with exact answers and
a uniform, algebraically derived charged-trace cost bound; an explicit
execution model with erasure, footprint-determinacy, payload-backing,
provenance, and mutation-rejection theorems; and an information-theoretic
lower bound in matching integer arithmetic. The one result held apart ---
an allocated-capacity theorem over a single packed object queried by a
closed controller --- has since been accepted in the repository, and this
draft will absorb it at the marked insertion point, and only there. What
remains genuinely open is listed in Section~\ref{sec:limitations}: the
controller's uncharged steps, preprocessing complexity, tightness of both
the constant and the overhead envelope, and mechanized cell-probe
redundancy bounds.

\bibliographystyle{alpha}
\bibliography{references}

\end{document}
